\section{Описание процесса изготовления продукта}

На данном этапе организационно-экономического обоснования определяется
состав и последовательность работ, необходимых для разработки методики
верификации соответствия TLA+ спецификации и исходного кода, а также
программного инструмента, реализующего предложенный подход.

Разработка проекта включает ряд этапов, начиная с анализа предметной
области и заканчивая подготовкой документации. Декомпозиция проекта
позволяет определить структуру работ, оценить их трудоёмкость и
сформировать основу для дальнейшего расчёта затрат.

\begin{table}[h]
\centering
\caption{Этапы разработки проекта}
\label{tab:stages}
\begin{tabular}{|c|p{10cm}|}
\hline
Этап & Наименование этапа \\ \hline
1 & Анализ предметной области и формирование требований \\ \hline
2 & Проектирование архитектуры системы \\ \hline
3 & Разработка механизма трассировки \\ \hline
4 & Разработка системы обработки трасс в TLA+ \\ \hline
5 & Тестирование и отладка \\ \hline
6 & Подготовка документации \\ \hline
\end{tabular}
\end{table}

Перечень событий и работ проекта приведён в таблице~\ref{tab:works}.

\begin{table}[h]
\centering
\caption{Перечень событий и работ проекта}
\label{tab:works}
\resizebox{\textwidth}{!}{
\begin{tabular}{|c|p{3cm}|p{7cm}|c|c|}
\hline
\parbox{2.5cm}{\centering № события\\ или работы} &
Событие & Работа & чел.-часы & чел.-дни \\ \hline

1 & Начало работ & -- & 0 & 0 \\ \hline

2 & Этап 1 &
Анализ предметной области и постановка задачи верификации &
32 & 4 \\ \hline

3 & &
Исследование языка TLA+, TLC и существующих подходов к
trace-checking &
40 & 5 \\ \hline

4 & &
Формирование функциональных и нефункциональных требований к
системе &
24 & 3 \\ \hline

5 & Этап 2 &
Проектирование архитектуры программного инструмента &
40 & 5 \\ \hline

6 & &
Разработка структуры модулей и интерфейсов взаимодействия &
24 & 3 \\ \hline

7 & &
Выбор технологий реализации и форматов хранения трасс &
20 & 2.5 \\ \hline

8 & Этап 3 &
Разработка библиотеки инструментирования исходного кода &
112 & 14 \\ \hline

9 & Этап 4 &
Разработка подсистемы агрегации трасс в единую трассу &
24 & 3 \\ \hline

10 & &
Разработка подсистемы обработки трасс в TLA+ и сравнение со спецификацией &
56 & 7 \\ \hline

11 & Этап 5 &
Тестирование, отладка и исправление выявленных ошибок &
44 & 5.5 \\ \hline

12 & Этап 6 &
Подготовка технической документации и описания методики &
20 & 2.5 \\ \hline

13 & Окончание работ & -- & 0 & 0 \\ \hline

\end{tabular}
}
\end{table}

Оценка трудоёмкости выполняется на основе экспертных оценок.
Для каждой работы задаются минимальное и максимальное значения
времени выполнения, после чего рассчитывается ожидаемая
трудоёмкость:

\begin{equation}
T = \frac{3T_{\min} + 2T_{\max}}{5}
\end{equation}

где $T_{\min}$ --- минимальная оценка времени выполнения работы,
$T_{\max}$ --- максимальная оценка.

Результаты расчёта трудоёмкости приведены в
таблице~\ref{tab:labor}.

\begin{table}[h]
\centering
\caption{Трудоёмкость работ проекта}
\label{tab:labor}
\resizebox{\textwidth}{!}{
\begin{tabular}{|c|c|p{6.6cm}|c|c|c|c|c|}
\hline
Этап & № & Содержание работы & $T_{\min}$ & $T_{\max}$ & $T$ &
чел.-дни & Исполнитель \\ \hline

1 & 1 &
Анализ предметной области и постановка задачи верификации &
28 & 38 & 32 & 4 & Аналитик \\ \hline

1 & 2 &
Исследование языка TLA+, TLC и существующих подходов к
trace-checking &
36 & 46 & 40 & 5 & Аналитик \\ \hline

1 & 3 &
Формирование функциональных и нефункциональных требований к
системе &
20 & 30 & 24 & 3 & Аналитик \\ \hline

2 & 4 &
Проектирование архитектуры программного инструмента &
36 & 46 & 40 & 5 & Архитектор \\ \hline

2 & 5 &
Разработка структуры модулей и интерфейсов взаимодействия &
20 & 30 & 24 & 3 & Архитектор \\ \hline

2 & 6 &
Выбор технологий реализации и форматов хранения трасс &
16 & 26 & 20 & 2.5 & Архитектор \\ \hline

3 & 7 &
Разработка библиотеки инструментирования исходного кода &
104 & 124 & 112 & 14 & Разработчик \\ \hline

4 & 8 &
Разработка подсистемы агрегации трасс в единую трассу &
20 & 30 & 24 & 3 & Разработчик \\ \hline

4 & 9 &
Разработка подсистемы обработки трасс в TLA+ и сравнение со спецификацией &
52 & 62 & 56 & 7 & Разработчик \\ \hline

5 & 10 &
Тестирование, отладка и исправление выявленных ошибок &
40 & 50 & 44 & 5.5 & Разработчик \\ \hline

6 & 11 &
Подготовка технической документации и описания методики &
16 & 26 & 20 & 2.5 & Аналитик \\ \hline

\multicolumn{5}{|r|}{Итого} & 436 & 54.5 & \\ \hline

\end{tabular}
}
\end{table}

Таким образом, суммарная трудоёмкость разработки проекта составляет
436 человеко-часов, что соответствует 54.5 человеко-дням. Наиболее
трудоёмким этапом является разработка механизма трассировки, поскольку
именно на этом этапе реализуется библиотека инструментирования
исходного кода, обеспечивающая фиксацию событий исполнения и
изменений переменных модели.

Для более детального анализа структуры трудоёмкости программных работ
используется разложение общей трудоёмкости каждой работы на
составляющие. Суммарные трудозатраты на выполнение $i$-й работы
представляются в виде:

\begin{equation}
q_i = q_i^{prog} + q_i^{al} + q_i^{test} + q_i^{doc},
\end{equation}

где $q_i^{prog}$ --- затраты труда на непосредственную реализацию
программного решения; $q_i^{al}$ --- затраты труда на алгоритмизацию;
$q_i^{test}$ --- затраты на тестирование и исправление ошибок;
$q_i^{doc}$ --- затраты на подготовку документации.

Затраты на алгоритмизацию определяются через коэффициент $n_a$:

\begin{equation}
q_i^{al} = n_a \cdot q_i^{prog}.
\end{equation}

Затраты на тестирование и исправление ошибок определяются по формуле:

\begin{equation}
q_i^{test} = q_i^{prog} \cdot (n_t + n_{cor}),
\end{equation}

где $n_t$ --- коэффициент тестирования, $n_{cor}$ --- коэффициент
внесения исправлений.

Затраты на подготовку документации:

\begin{equation}
q_i^{doc} = n_d \cdot q_i^{prog}.
\end{equation}

Тогда суммарная трудоёмкость может быть выражена как:

\begin{equation}
q_i = q_i^{prog}(1 + n_a + n_t + n_{cor} + n_d).
\end{equation}

Примем значения коэффициентов, рекомендуемые методикой:

\[
n_a = 0.3, \quad n_t = 0.3, \quad n_{cor} = 0.3, \quad n_d = 0.35.
\]

Тогда:

\[
q_i = q_i^{prog}(1 + n_a + n_t + n_{cor} + n_d) =
q_i^{prog} \cdot 2.25.
\]

Отсюда:

\begin{equation}
q_i^{prog} = \frac{q_i}{2.25}.
\end{equation}

Затраты на отдельные составляющие определяются как:

\[
q_i^{al} = 0.3q_i^{prog}, \quad
q_i^{test} = 0.6q_i^{prog}, \quad
q_i^{doc} = 0.35q_i^{prog}.
\]

Рассмотрим пример расчёта структуры трудоёмкости для этапа 1
«Анализ предметной области и формирование требований».

Согласно таблице~\ref{tab:labor}, суммарная трудоёмкость этапа
составляет:

\[
q_1 = 32 + 40 + 24 = 96 \text{ чел.-ч}.
\]

Определим трудозатраты на программную реализацию:

\[
q_1^{prog} = \frac{q_1}{2.25} = \frac{96}{2.25} \approx
42.67 \text{ чел.-ч}.
\]

Тогда остальные составляющие равны:

\[
q_1^{al} = 0.3 \cdot 42.67 \approx 12.80 \text{ чел.-ч},
\]

\[
q_1^{test} = 0.6 \cdot 42.67 \approx 25.60 \text{ чел.-ч},
\]

\[
q_1^{doc} = 0.35 \cdot 42.67 \approx 14.93 \text{ чел.-ч}.
\]

Аналогичным образом выполняются расчёты для остальных этапов.
Результаты приведены в таблице~\ref{tab:structure}.

\begin{table}[h]
\centering
\caption{Структура трудоёмкости работ по этапам проекта}
\label{tab:structure}
\resizebox{\textwidth}{!}{
\begin{tabular}{|c|p{6cm}|c|c|c|c|c|}
\hline
Этап & Наименование этапа & $q_i$, чел.-ч &
$q_i^{prog}$ & $q_i^{al}$ & $q_i^{test}$ & $q_i^{doc}$ \\ \hline

1 & Анализ предметной области и формирование требований &
96 & 42.67 & 12.80 & 25.60 & 14.93 \\ \hline

2 & Проектирование архитектуры системы &
84 & 37.33 & 11.20 & 22.40 & 13.07 \\ \hline

3 & Разработка механизма трассировки &
112 & 49.78 & 14.93 & 29.87 & 17.42 \\ \hline

4 & Разработка системы обработки трасс в TLA+ &
80 & 35.56 & 10.67 & 21.33 & 12.44 \\ \hline

5 & Тестирование и отладка &
44 & 19.56 & 5.87 & 11.73 & 6.84 \\ \hline

6 & Подготовка документации &
20 & 8.89 & 2.67 & 5.33 & 3.11 \\ \hline

\multicolumn{2}{|r|}{Итого} & 436 & 193.79 & 58.14 & 116.26 & 67.81 \\ \hline

\end{tabular}
}
\end{table}
